\documentclass[10pt,twocolumn]{article}
\usepackage[a4paper,margin=15mm]{geometry}
\usepackage{amsmath,amssymb}
\usepackage{mathrsfs}
\usepackage{amsthm}
\theoremstyle{definition}
\newtheorem{theorem}{Theorem}
\newtheorem{proposition}{Proposition}
\newtheorem{corollary}{Corollary}
\newtheorem{remark}{Remark}
\newtheorem{definition}{Definition}
\usepackage{graphicx}
\usepackage{booktabs,array,calc}
\usepackage[font=small,labelfont=bf]{caption}
\usepackage[colorlinks=true]{hyperref}
\usepackage{etoolbox}
\providecommand{\tightlist}{\setlength{\itemsep}{0pt}\setlength{\parskip}{0pt}}
\setcounter{secnumdepth}{-1}
\emergencystretch=3em
\allowdisplaybreaks
\AtBeginDocument{%
  \BeforeBeginEnvironment{equation}{\small}%
  \BeforeBeginEnvironment{equation*}{\small}%
  \BeforeBeginEnvironment{align}{\small}%
  \BeforeBeginEnvironment{align*}{\small}%
  \BeforeBeginEnvironment{gather}{\small}%
  \BeforeBeginEnvironment{gather*}{\small}%
}

\begin{document}
\title{Vector Floors: Farkas-Certified Obstruction\\ under Multiple Sustainability Constraints\\[0.2em]
{\large\rmfamily Edition 2, revised after a joint audit}}
\author{Amin Abaee\\[0.35em]
{\small Independent Researcher}\\[0.55em]
{\small\href{https://orcid.org/0000-0002-0019-1842}{ORCID: 0000-0002-0019-1842}}\\[0.3em]
{\small\href{mailto:amin\_abaee@ut.ac.ir}{amin\_abaee@ut.ac.ir}}}
\date{September 24, 2026}
\maketitle

\begin{abstract}
Sustainability constraints are rarely single floors: minimum stocks, maximum
emissions, service levels, and institutional limits bind simultaneously. The calculus states its certificates
for general vector constraint sets $\mathcal V=\{x: q_j(x)\ge 0,\
j=1,\dots,m\}$ (Abaee, 2026, An obstruction calculus, Section 2.2), but its worked certificates and
audited instances are carried by a single scalar margin (Section 3.4's
two-floor certificate there). This paper, the D4 item of the programme's
generalization catalogue, works the vector layer as such: the stacked
safe-action system is infeasible precisely when a Farkas multiplier vector
exists, so the vector obstruction itself is a checkable object of rational
arithmetic; the new content over the source paper is the dual-certificate
reading, the floor-monotonicity count on the audited pair family, and the
three-floor support instance. On the programme's two-patch working system
(Abaee, 2026, An obstruction calculus, Section 8) the dynamic vector-floor reading reproduces the
audited structure verdicts (blind and aggregate readings nonviable;
$z_1$-only, $z_2$-only, and full readings viable), and adding a floor
strictly shrinks the viable pair family --- from $36$ of $36$ viable pairs
on one floor to $28$ on two, the pairs lost being exactly those containing
the dynamically dead corner $(1,1)$. A three-floor stacked instance
exhibits the support reading of the multipliers (the certificate's carrier
is a minimal infeasible subsystem: inactive row, multiplier zero). All
claims are machine-verified in exact rational arithmetic; the verification
script ships with the paper.
\end{abstract}

\section{Introduction and relation to the programme}

The programme's completed items --- the obstruction-calculus manuscript with
its four certificate families, the continuous-to-finite companion bridge,
the viable-selector unification, the stochastic selector, and the five-layer
successor absorbing general information structures and exit-time values ---
define vector
constraint sets from the outset ($\mathcal V=\{x: q_j(x)\ge0,\ j=1,\dots,m\}$,
Abaee, 2026, An obstruction calculus, Section 2.2), and the polyhedral Farkas equivalence for the
stacked safe-action system is already stated there (Section 3.2), with a
worked two-floor certificate (Section 3.4). What the completed items do not
do is carry the vector layer as an object of computation: the audited
instances are scalar-margin instances, and the catalogue's D4 proposal ---
real constraint systems are vectors of floors --- is left unexecuted. The
present paper executes exactly that: the vector common-action
correspondence, the stacked certificate as a checkable dual object, the
floor-monotonicity count on the audited pair family, and the three-floor
support instance. Nothing in the completed items is modified; the additions
are computational, and each is positioned against its source passage.

\section{The vector-floor model}

Let the state space carry $m$ floors $q_1,\dots,q_m$ with safe set
\begin{equation*}
\mathcal V=\{x:\ q_j(x)\ge 0\ \text{for every } j\}.
\end{equation*}
For a state $x$ with instrument set $U(x)$ and discrete-time transition
$x^{+}=F(x,u,d)$, the \emph{vector safe-action set} collects the instruments that keep every floor non-negative under every
admissible disturbance:
\begin{equation*}
\mathcal R_{\mathcal V}(x)=\{u\in U(x):\ x^{+}\in\mathcal V\
\text{for all successors } x^{+}\}.
\end{equation*}
For a belief $B$ (the set of states compatible with the reading), the
common-action condition with per-cell semantics reads
$\mathcal R^B_{\mathcal V}(B)=\bigcap_{x\in B}\mathcal R_{\mathcal V}(x)$
within each indistinguishability cell, exactly as in the successor's
read-then-act convention. The vector obstruction is the static statement
\begin{equation*}
\mathcal R^B_{\mathcal V}(B)=\varnothing\ \Longrightarrow\
B\notin\mathrm{ERViab}_{\mathcal I}(\mathcal V),
\end{equation*}
where $\mathrm{ERViab}_{\mathcal I}(\mathcal V)$ is the observation-based
viability kernel under structure $\mathcal I$ in the programme's notation
(Abaee, 2026, An obstruction calculus, Section 2.4): a single moment without a jointly
floor-preserving instrument certifies nonviability, no matter how capable
the future.

\begin{proposition}[Branchwise/joint gap for vector floors]
\label{prop:gap}
(The vector instance of the common-action obstruction, Abaee, 2026, An obstruction calculus,
Theorem 2; the vector-specific content is the stacked dual layer of
Proposition~\ref{prop:farkas} and the floor-monotonicity of E4.)
If for each state $x\in B$ separately $\mathcal R_{\mathcal V}(x)\neq
\varnothing$ but $\bigcap_{x\in B}\mathcal R_{\mathcal V}(x)=\varnothing$,
then every branch of $B$ is individually floor-servable while the belief is
vector-obstructed: the information constraint, not any single floor, is the
cause of nonviability.
\end{proposition}
\begin{proof}
The branchwise hypothesis provides, for each $x$, an instrument $u_x$
keeping $x$ in $\mathcal V$ under \emph{all} floors jointly; each branch
is thus servable in isolation. The joint hypothesis denies the existence of
one instrument admissible for all $x$ simultaneously, which is precisely
the per-cell read-then-act requirement; any observation-based policy must
move the whole cell with one instrument and therefore fails at the first
step. The cause cannot be attributed to any single floor's infeasibility
(every branch is jointly floor-servable), nor to the future (the failure
occurs in one step): it is the pairing of one instrument with several
states, i.e.\ the information constraint.
\end{proof}

\section{The Farkas certificate}

Let the floors and the instrument set be polyhedral, so that each
safe-action set is $\mathcal R_{\mathcal V}(x_i)=\{u: A_i u\le b_i\}$ with
rational data. The common safe-action system stacks to $Au\le b$, and:

\begin{proposition}[Obstruction as dual infeasibility]
\label{prop:farkas}
The stacked system $Au\le b$ is infeasible if and only if there exists a
rational $\lambda\ge 0$ with $\lambda^{\top}A=0$ and $\lambda^{\top}b<0$
(Farkas' lemma). Consequently the vector obstruction admits a certificate
that is itself checkable in exact rational arithmetic: given $\lambda$, the
three identities are verifiable by Fraction arithmetic with no solver and
no search.
\end{proposition}
\begin{proof}
Farkas' lemma gives the equivalence. Checkability: the certificate side
requires evaluating two matrix products and one inequality --- exact
operations on rationals. The infeasibility side is certified, not searched:
$\lambda^{\top}A=0$ and $\lambda^{\top}b<0$ with $\lambda\ge0$ yield
$0=\lambda^{\top}(Au)\le\lambda^{\top}b<0$ for every candidate $u$, a
contradiction.
\end{proof}

\begin{remark}[Support reading]
In the stacked system of a three-floor instance below, the certified
multiplier assigns weight $0$ to a redundant row: the certificate's
\emph{support} is a minimal infeasible subsystem of the stacked rows (an
infeasible system admits no primal solution, so the reading is one of dual
support, not of primal activity). The obstruction is carried by exactly
the floors that bind.
\end{remark}

\section{Verified instances}

All of the following are re-derived and checked by the verification script
(\textsf{8/8 exact check families}; standard library, Fraction arithmetic,
deterministic).

\paragraph{E1 (stacked equivalence, both directions).}
The two-floor instrument instance $\mathcal R_1=\{u: u\le\tfrac25\}$,
$\mathcal R_2=\{u: u\ge\tfrac35\}$ stacks to rows $u\le\tfrac25$,
$-u\le-\tfrac35$; the multiplier $\lambda=(1,1)$ satisfies
$\lambda^{\top}A=0$ and $\lambda^{\top}b=\tfrac25-\tfrac35=-\tfrac15<0$,
certifying infeasibility; a $101$-point rational grid confirms no $u$
satisfies both. The feasible contrast $\{u\le\tfrac25\}\cap\{u\ge\tfrac15\}$
(with a redundant third row $u\le\tfrac45$) admits no certificate --- the
exact case analysis forces $\lambda_1=\lambda_2$ and then
$\lambda^{\top}b\ge 0$ --- and the witness $u=\tfrac15$ is exhibited. Both
directions of Proposition~\ref{prop:farkas} are thereby exercised; the
instance is the unnormalized form of the worked certificate of Abaee, 2026, An obstruction calculus,
Section 3.4 (there with the multiplier normalized to $\|\lambda\|_1=1$),
and the feasible contrast is the direction that source states but does not
exercise.

\paragraph{E2 (the belief-level gap).}
On the belief merging a state with floor-set $\{u\le\tfrac25\}$ with a
state with floor-set $\{u\ge\tfrac35\}$: branchwise witnesses $u=\tfrac2{10}$
and $u=\tfrac6{10}$ exist; the joint set is empty with the E1 certificate.
Proposition~\ref{prop:gap} fires exactly.

\paragraph{E3 (dynamic vector floors).}
On the programme's two-patch system --- patches $i\in\{1,2\}$ with stocks
$z_i\in\{0,1,2,3\}$, regeneration $+1$ per step capped at $3$, damage $2$
on the unprotected patch, instruments $u\in\{1,2\}$ protecting one patch;
defined in Abaee, 2026, An obstruction calculus, Section 8, and restated with full transition data
in the verification script --- the safe set is already a vector object
$q_1: z_1\ge 1$, $q_2: z_2\ge 1$ ($m=2$) served by one shared instrument.
Per-cell read-then-act viability over the audited beliefs
$\{(1,2),(2,1)\}$ and $\{(1,2),(2,1),(2,2)\}$ reproduces the audited
verdicts under five information structures: blind nonviable, aggregate
nonviable, $z_1$-only viable, $z_2$-only viable, full viable --- the
successor's structure table re-read as a statement about two floors and one
instrument.

\paragraph{E4 (floor monotonicity).}
The \emph{viable pair family}: the $\binom{9}{2}=36$ two-state beliefs of
the nine floor-safe states, a pair being viable when both of its states are
(full observation splits the pair into its branches). Adding the second
floor strictly shrinks the family from all $36$ to $28$: under two floors
the corner $(1,1)$ has \emph{no} instrument keeping both floors (one
shared instrument, damage $2$ on the other patch), so $(1,1)$ is
dynamically nonviable and exactly the $8$ pairs containing it are lost.
Dropping the second floor ($q_1$ only) re-admits all $36$, setwise, with
strict separation witnessed by $\{(1,1),(1,2)\}$: viable on one floor,
nonviable on two. Adding floors destroys viability --- the marginal cost of
each additional institutional constraint is a computable set.

\paragraph{E5 (exact multiplier recheck).}
Every multiplier used above is re-validated componentwise
($\lambda\ge0$; $\lambda^{\top}A=0$; $\lambda^{\top}b<0$) in Fraction
arithmetic.

\paragraph{E6 (three floors, complementarity).}
The stacked system $\{u\le\tfrac25\}\wedge\{u\ge\tfrac35\}\wedge
\{u\le\tfrac9{10}\}$ is infeasible with certificate $\lambda=(1,1,0)$: the
inactive third row carries multiplier zero, exhibiting the support reading
of Remark~1 (the script further certifies minimality: every two-row
subsystem containing row 2 or row 3 is feasible).

\section{What is not claimed}

The dynamic vector-floor theory is not developed beyond the per-cell
re-reading of the audited systems: no vector-valued comparison functions,
no multi-floor dynamic programming, no Pareto readings of trade-offs
between floors. The Farkas layer is static; the companion bridge's
continuous-to-finite transfer and the successor's timing layer are cited,
not re-derived. Non-convex instrument codices invalidate
Proposition~\ref{prop:farkas} (Farkas' lemma is a convexity statement);
the monitoring-design companion of this catalogue records the exact
non-convex counterinstance. Sustainability floors with uncertain
thresholds are treated in the programme's buffered-viability companion, not
here. The dynamic vector layer is likewise recorded, not developed: the
first-breach timing bound for floors with margins $\varepsilon_j$ --- exit
by $\min_j \inf_{\mathcal V} q_j/\varepsilon_j$ --- and the multi-face
Isaacs reading at corners where several floors are simultaneously active
belong to the successor's timing layer; so does the distinction between
cross-floor infeasibility at one state and cross-state infeasibility under
merged readings, which Proposition~\ref{prop:gap} isolates in its static,
one-step form.

\section*{Declarations}
\textbf{Funding.} None.\quad \textbf{Competing interests.} None.\quad
\textbf{Data availability.} No external data were used.\quad
\textbf{Code availability.} The verification script
\texttt{vector\_floor\_certificates\_v2\_verify.py} is publicly available at
\url{https://github.com/MIKEAA2020/general-sustainability} (folder
\texttt{arena agent 1/paper rewrites/latex}); it reproduces all eight check
families verbatim.\quad
\textbf{AI declaration.} GLM (Z.ai), Qwen (Alibaba Cloud) and DeepSeek AI
assisted with drafting and iterative review of the first edition; the joint
audit of the first edition and this revision were prepared with further AI
assistance, under the programme's verification discipline.

\subsection{References}\label{references}
{\footnotesize

Abaee, A.: An obstruction calculus for viability under incomplete
observation. Submitted.
\url{https://github.com/MIKEAA2020/general-sustainability} (2026)

}

\end{document}
